\begin{musica}{Olhos Negros}{Popular Açoreana}{Tuna Universitária do Porto}
\nota{Introdução: G}

Os teus \ac{G}{o}lhos negros, negros,
São gen\ac{D}{ti}os, são gentios da Gui\ac{G}{né}.
Os teus olhos negros, negros
São gen\ac{D}{ti}os, são gentios da Gui\ac{G}{né}.

\ac{g}{Ai}, \ac{a}{da} \ac{b}{Gu}i\ac{C}{né} \ac{Cm}{por} serem \ac{G}{ne}gros,
\ac{g}{Da} \ac{f#}{Gu}i\ac{Em}{né} por serem ne\ac{Am}{gros},
Gen\ac{D7}{tios} por não terem \ac{G}{fé}.

Olhos negros são brilhantes,
Semelhantes aos luzeiros que o céu tem.
Olhos negros são brilhantes,
Semelhantes aos luzeiros que o céu tem.

Ai, eu amei dois olhos negros,
Eu amei dois olhos negros
Sem fazer mal a ninguém.

Os teus olhos negros, negros,
São a noite e são estrelas que alumiam.
Os teus olhos negros, negros,
São a noite e são estrelas que alumiam.

São a noite porque me perdem,
São a noite porque me perdem,
São estrelas porque me guiam.

Uns olhos negros eu preferi,
Eu preferi e nunca vi
De cor linda em ninguém.

\posicoes{D: xx0232, D7: xx0212, Em: 022000, G: 320033, Am: x02210}
\end{musica}
